\chapter{Invariants and monovariants}

\begin{orient}
\textbf{What this chapter is for.} This is the highest-yield idea in combinatorial problem solving, and it deserves the claim. Whenever a problem describes a process, a move, an operation repeated any number of times, there is a question worth asking before anything else: what does the move fail to change? If some quantity is preserved, then every reachable state agrees with the start on that quantity, and any target that disagrees is unreachable. A search that could run forever collapses to one line of arithmetic. The companion idea, the monovariant, is a quantity that moves only one way; it cannot prove unreachability but it proves termination, and it bounds how long a process can run. The difficult part is never using an invariant. It is finding one, and this chapter is organised around that search.
\end{orient}

\section{Finding the invariant}

There is a reliable order in which to look. Parity first, because it is free and it works surprisingly often. Then a colouring, because a colouring is a parity argument with more colours and it is the natural tool when the moves are geometric. Then a weighting, when the objects have positions or sizes that the move relates. Then an algebraic combination, typically a sum or a product taken modulo a small number. If all four fail, the problem is probably not an invariant problem, and the chapter on counting or the chapter on games is the right place to look next.

\tech{Parity invariant}{easy}
\cue A move that changes several quantities at once by amounts summing to an even number, or that flips signs, or that swaps two objects. Any question of the form ``can this configuration become that one'', where the two configurations look similar.
\method Find a quantity whose parity the move preserves: a count of objects of one kind, a sum, the number of inversions in an arrangement, the number of objects in a distinguished region. Then compare the parity of the start with the parity of the target. Different parity means unreachable, and the proof is complete.

\begin{worked}
A board carries the numbers $1$ through $2n$. A move deletes two numbers $a,b$ and writes $\abs{a-b}$. After $2n-1$ moves one number remains. Is it even or odd? The move replaces $a+b$ by $\abs{a-b}$, and $a+b$ and $\abs{a-b}$ always have the same parity, so the parity of the total sum is invariant. The initial sum is $1+2+\cdots+2n=n(2n+1)$, whose parity is that of $n$. So the final number is even exactly when $n$ is even. For $n=3$ the initial sum is $21$, odd, so the survivor is odd whatever moves are made, which is easy to confirm by random simulation.
\end{worked}

\begin{trap}
Verifying the invariant on one move and assuming it holds for all. Check every move type the problem allows, including the degenerate ones, and check what happens at the boundary of the board or the end of the list.
\end{trap}

\begin{seealsob}Colouring arguments; multiplicative and other algebraic invariants; tiling impossibility by colouring.\end{seealsob}

\tech{Colouring arguments: chessboard and modulo $k$}{medium}
\cue A geometric or grid problem in which a piece covers a fixed shape, or moves by a fixed step, and you are asked whether a covering or a tour is possible. Also any tiling question.
\method Colour the cells so that every legal placement or every legal move covers a fixed multiset of colours. Then count colours in the region and compare. Two colours in a chequered pattern handle dominoes and rook or bishop moves; colouring by column index modulo $k$ handles $1\times k$ tiles; diagonal colourings handle knight-like steps. The colouring is chosen to make the move's colour signature constant, and inventing it is the entire creative act.

\begin{worked}
Remove two opposite corners from an $8\times8$ board and try to tile the remaining $62$ squares with $31$ dominoes. Colour the board normally. Every domino covers one black and one white square, so a tiling needs equal counts. Opposite corners share a colour, so the mutilated board has $32$ of one colour and $30$ of the other. No tiling exists.
\end{worked}

\begin{worked}[A colouring with more than two colours]
Can a $10\times10$ board be tiled by $1\times4$ pieces? Colour cell $(i,j)$ with $(i+2j)\bmod 4$. Every horizontal piece covers the four residues once each; every vertical piece covers $(i+2j)$, $(i+1+2j)$, $(i+2+2j)$, $(i+3+2j)$, again all four residues once each. So a tiling requires the four colour classes to be equal in size, at $25$ each. Counting the classes on the $10\times10$ board gives $25,25,25,25$, so this colouring does not obstruct. The standard four-colouring by $(i+j)\bmod 4$ does not obstruct either, and in fact a tiling exists, which is the honest outcome: a colouring that balances proves nothing, and you must then construct.
\end{worked}

\begin{trap}
Concluding possibility from a colouring that balances. A colouring argument is a one-way test: imbalance proves impossibility, balance proves nothing at all. When the count comes out even, you still owe a construction.
\end{trap}

\begin{seealsob}Parity invariant; weighting arguments; tiling impossibility by colouring.\end{seealsob}

\tech{Weighting arguments}{hard}
\cue A move that relates positions with different values, most characteristically a jump in which a piece at position $p$ moves to position $p+2$ while a piece at $p+1$ is removed. Also any problem asking how far something can reach.
\method Assign a weight $w(c)$ to each cell so that the move never increases the total weight, then compare the target's weight with the finite total available. The classic choice is a geometric weighting $w=\varphi^{d}$ where $d$ is the distance to the target and $\varphi$ is the golden ratio, chosen exactly because $\varphi^{2}=\varphi+1$ makes a jump weight-preserving.

\begin{worked}
Conway's soldiers. Pegs fill the half-plane $y\leq0$ of a grid, and a move jumps a peg over an adjacent peg into an empty cell two away, removing the jumped peg. How far above the line can a peg reach? Weight cell $c$ by $\varphi^{d(c)}$ where $d$ is the taxicab distance to the target cell and $\varphi=(\sqrt5-1)/2$ satisfies $\varphi^{2}+\varphi=1$. A jump toward the target preserves weight, and any other jump strictly decreases it. Summing the weights of the entire filled half-plane converges, and the total is exactly $1$, the weight of a peg sitting on the target itself. Reaching level $5$ would require a weight strictly greater than the whole half-plane supplies, so level $5$ is unreachable no matter how many pegs are used.
\end{worked}

\begin{trap}
Choosing a weighting that is preserved by the move but whose total over the initial configuration diverges. The argument needs a finite bound to compare against, so the geometric decay is not decoration.
\end{trap}

\begin{seealsob}Potential functions; monovariants; colouring arguments.\end{seealsob}

\tech{Multiplicative and other algebraic invariants}{medium}
\cue Moves that multiply, that change signs, or that act on a tuple by a linear substitution. Also processes on numbers where the sum is visibly not invariant.
\method Try the product of the entries, the product of the signs, the sum of squares, the greatest common divisor, or a linear combination $\sum a_{i}x_{i}$ with coefficients chosen so that the move's effect cancels. The last of these is a small linear algebra problem: write the move as a vector of changes and find a functional that annihilates it.

\begin{worked}
Numbers $a,b,c$ on a board; a move replaces them by $a+b-c$, $b+c-a$, $c+a-b$ in some order. The sum $a+b+c$ is invariant, since the three new numbers sum to $a+b+c$. So starting from $(1,2,3)$ you can never reach $(2,3,5)$, whose sum is $10$ rather than $6$, whatever sequence of moves is used.
\end{worked}

\begin{trap}
Finding an invariant and forgetting to check that the target genuinely violates it. If start and target agree on every invariant you found, you have learned nothing and must either find another invariant or construct a path.
\end{trap}

\begin{seealsob}Parity invariant; weighting arguments; residue classes and modular reduction.\end{seealsob}

\section{Quantities that move one way}

\tech{Monovariants and the termination argument}{medium}
\cue The question is not ``can this be reached'' but ``must this stop'', ``how many moves at most'', or ``what does it look like when no move is available''.
\method Find a quantity that strictly decreases (or strictly increases) with every legal move and that is bounded, and preferably integer-valued. A strictly decreasing non-negative integer cannot decrease forever, so the process terminates, and the bound on the quantity bounds the number of moves. Then characterise the terminal states: they are exactly the configurations in which no move is legal, and identifying them is often the actual answer.

\begin{worked}
People sit around a table; if any two adjacent people are ordered wrongly by height, they may swap. Does this stop? The number of inversions, meaning pairs out of order, strictly decreases by one at every legal swap, and is a non-negative integer bounded by $\binom{n}{2}$. So the process halts after at most $\binom{n}{2}$ swaps, and it halts only when no adjacent pair is inverted, which for a linear order means fully sorted.
\end{worked}

\begin{trap}
A monovariant that decreases weakly rather than strictly. ``Does not increase'' permits an infinite loop among states of equal value, so either sharpen the quantity or add a tie-breaking second component and argue lexicographically.
\end{trap}

\begin{seealsob}Potential functions; well-ordering and the minimal counterexample; winning and losing positions.\end{seealsob}

\tech{Potential functions}{hard}
\cue A process where no single obvious quantity is monotone, but a weighted combination might be. Also amortised bounds, where individual moves can be expensive but the total is not.
\method Design $\Phi$ on states so that every move decreases $\Phi$ by at least a fixed positive amount, or so that the actual cost of a move plus the change in $\Phi$ is bounded. Then the total number of moves is at most $\left(\Phi_{\text{start}}-\Phi_{\min}\right)/\delta$. Designing $\Phi$ is the same craft as designing a weighting, and the usual candidates are sums of squares, sums of a convex function of the parts, and entropy-like expressions.

\begin{worked}
Chip-firing on a line of $n$ cells: a cell holding at least two chips may send one chip to each neighbour. Take $\Phi=\sum_{i}c_{i}\,i^{2}$ where $c_{i}$ is the chip count at cell $i$. Firing at $i$ changes $\Phi$ by $(i-1)^{2}+(i+1)^{2}-2i^{2}=2$, so $\Phi$ strictly increases by exactly $2$ per firing while remaining bounded above by $n^{2}$ times the total chip count. The process therefore terminates, and the number of firings is determined by the change in $\Phi$ alone, independent of the order in which cells are fired.
\end{worked}

\begin{seealsob}Monovariants; weighting arguments; the extremal principle.\end{seealsob}

\section{Recognition matrix}

\begin{recog}{@{}p{0.31\textwidth}p{0.35\textwidth}p{0.28\textwidth}@{}}
\rowcolor{mist}\mh{If you see} & \mh{Reach for} & \mh{If that fails} \\
\midrule
\endhead
``Can state $X$ become state $Y$'' & An invariant separating them & Construct a path \\
A move changing two things at once & Parity of the sum & Sum modulo a small $k$ \\
Tiles on a grid & Colouring matched to the tile & Weighted colouring \\
$1\times k$ tiles & Colour by index modulo $k$ & Diagonal colouring \\
Jumps of length two & Geometric weighting with $\varphi^{2}=\varphi+1$ & Sum of squares \\
Signs flipping & Product of the signs & Count of negatives modulo 2 \\
Linear substitution on a tuple & A functional annihilating the change & Eigenvector of the move matrix \\
``Show the process stops'' & A strictly decreasing integer & Lexicographic pair \\
``How many moves at most'' & Bound the monovariant's range & Potential function \\
Order matters, apparently & Check whether $\Phi$ change is order-free & Abelian property \\
Colouring balances & Nothing proved; construct & Try a finer colouring \\
\end{recog}
